\section{Recovering coordinate reflection symmetry}\label{sec:coordinate-symmetry}

The local operator estimate first gives reflection symmetry for
observables on the two layers adjacent to a coordinate link plane.
To use this estimate for an arbitrary curvature character, we remove
its outer fermion fields by exact integration identities and its
outer photon exponents by a finite contour shift. Only finitely many
terms of the reflected perturbation change in each shift; their
contribution is of order its strength, uniformly in volume. We reproduce
the two-layer regularization of \cite[Sections~8.9--8.10]{ChamberQED}
and the elimination argument of its Section~9, culminating in
Theorem~9.25 there.

\subsection{The two adjacent fermion layers}\label{subsec:diag-two-layers}

The remaining comparison argument uses gauge invariant polynomials
\[
 P(\widehat A)=\sum_a b_a e^{i\langle\delta_a,\widehat A\rangle}
                                     \Psi_a,\qquad
 d^*\delta_a=q\mathfrak q(\Psi_a),
\]
where the cochains have finite support and the $\Psi_a$ are ordered
Grassmann monomials. Denote their class by $\mathcal P$. The charge
condition is exactly gauge invariance. It also makes each term
even: the total divergence is zero, so its numbers of $\psi$ and
$\bar\psi$ factors agree. The expansion is unique after equal
pairs $(\delta_a,\Psi_a)$ have been collected, because the monomials
and the distinct characters are linearly independent. Each such
polynomial is an admissible local integrand with bounded coefficients.

Write $v_0=(v_1,v_2,v_3)$, with $0<v_3<v_2<v_1<1/2$, and set
\[
 (s_1,\ldots,s_6)=(-v_1,-v_2,-v_3,v_3,v_2,v_1),\qquad
 \mathcal X_{j,a}=\{x\in\mathcal X:x_1=j+s_a\}.
\]
These are the six thin layers of the thick layer centered at $x_1=j$.
At a coordinate link plane $c=j+\tfrac12$, a two-layer polynomial
reads photon links in the planes $j$ and $j+1$ and between them,
and fermion fields only at the two ends of the bonds crossing
$c$. These are the thin layers $\mathcal X_{j,6}$ and
$\mathcal X_{j+1,1}$; the other thin layers of the two neighboring thick layers do
not occur. Call this subclass $\mathcal P_1(c)$.

Two estimates connect this class to coordinate reflection symmetry.
First, for a monomial $Y\in\mathcal P_1(c)$, heat regularization constructs a
bounded operator $X_{Y,\epsilon}$ and an error
$\xi_Y(\epsilon)\to0$ such that
\[
 |\mu_i^t(\Theta_cY)-\overline{\mu_i^t(Y)}|
 \le4\sqrt{\gamma_{t,i}(k)^2-1}\,\|X_{Y,\epsilon}\|
       +C_Y\gamma_{t,i}(k)\xi_Y(\epsilon).
\]
The fixed support scale $k$ is independent of $\epsilon$.
The operator norm can grow exponentially in $1/\epsilon$; thus the
volume limit must precede $t\downarrow0$, and only then may
$\epsilon\downarrow0$. The construction and the monomial form of
this estimate appear in \eqref{eq:diag-two-layer-finite-error};
finite sums give the symmetry assertion for every two-layer polynomial.
Second, \cref{thm:diag-polynomial-reduction} assigns to every
$P\in\mathcal P$ a $Y_P\in\mathcal P_1(c)$, independent of the
torus and tilt, with
\[
 |\mu_i^t(P^g)-\mu_i^t((Y_P)^g)|\le C_{P,g}t
 \qquad(0\le t\le1,\ i\ge i_0(P,g))
\]
for each fixed lattice isometry $g$. Its proof removes the outer
fermion layers and then shifts the photon contour. This estimate
extends the two-layer symmetry to the curvature characters used
to identify the limiting law.

For a monomial
$Y=e^{i\delta\widehat A}\Psi$, write its exponent as
$(\delta_0,\delta_1,\delta_\perp)$ on the lower spatial, upper
spatial and temporal links. If $\mathfrak q_\pm$ are the lower
and upper charge densities, gauge invariance gives
\begin{equation}\label{eq:diag-two-layer-gauss}
 d_s^*\delta_0-\delta_\perp=q\mathfrak q_-,\qquad
 d_s^*\delta_1+\delta_\perp=q\mathfrak q_+.
\end{equation}
Let $S=\supp\delta_0\cup\supp\delta_1$ and choose a nonempty $\Lambda$
containing their endpoints, the support of $\delta_\perp$ and
all sites over which the fields of $\Psi$ occur.

Every such $Y$ is a kernel insertion, but the operator between
the usual quarter-heat factors need not be bounded. We give the
different regularization that permits approximation by bounded
operators. The crossing modes and their reflected modes form a
basis of the two-layer generators. Since all their weights are
positive when $0<\rho<1$, there is a unique finite matrix
$\Upsilon_Y$ with
\begin{equation}\label{eq:diag-two-layer-matrix}
 \left.\Psi e^{-S_c^\times(\Lambda)}\right|_{\widehat a_j=0}
       =\operatorname{Ins}_c[
                      \omega^{1/2}\Upsilon_Y\omega^{1/2}].
\end{equation}
To obtain it, expand the left side in the basis
$\Theta_c\zeta_J\zeta_I$ and multiply the resulting matrix by
$\omega^{-1/2}$ on both sides. Its nonzero entries have
$\mathfrak q(I)-\mathfrak q(J)=\mathfrak q_-+\mathfrak q_+$.
Thus $\Upsilon_Y$ has charge
$q^{-1}d_s^*(\delta_0+\delta_1)$ and is even.

Set
\begin{equation}\label{eq:diag-two-layer-kernel}
 k_Y^0(A',A)=e^{i\delta_1A'}
             \prod_{b\in S}g_1(A'(b)-A(b))e^{i\delta_0A}
                         \Upsilon_Y.
\end{equation}
It is the continuous kernel of
$e^{i(\delta_0+\delta_1)A}e^{-|E_S+\delta_0|^2/2}\Upsilon_Y$.
We also write $k_Y^0$ for this integral operator.
The charge identity and unchanged Gaussian difference prove
its covariance. Moreover
\begin{equation}\label{eq:diag-two-layer-is-insertion}
 Y=G[k_Y^0].
\end{equation}
Indeed, the Gaussian factors in the insertion cancel. Applying
the automorphism that multiplies the lower fields by their
charge phases to \eqref{eq:diag-two-layer-matrix} replaces its
matrix by $\omega^{1/2}\Upsilon_Yu(\widehat a_j)\omega^{1/2}$
and its left side by
$e^{-iq\langle\mathfrak q_-,\widehat a_j\rangle}
 \Psi e^{-S_c^\times(\Lambda)}$. The remaining temporal exponent
in $G[k_Y^0]$ is
$d_s^*\delta_0-q\mathfrak q_-=\delta_\perp$ by
\eqref{eq:diag-two-layer-gauss}, proving the identity.

Put $\widetilde E=E_S+\delta_0$, $D=\delta_1-\delta_0$ and
$a=(|\delta_0|^2+|\delta_1|^2)/8$. Define the bounded operator
\begin{equation}\label{eq:diag-two-layer-prime}
 X'_Y=e^{i(\delta_0+\delta_1)A}
       e^{-\frac14|\widetilde E|^2+
                         \frac14\langle\widetilde E,D\rangle+a}
                                                      \Upsilon_Y.
\end{equation}
It belongs to $\mathfrak W(S,\Lambda)$, and completion of the
quadratic exponent gives
\[
 e^{-|E_S|^2/8}X'_Ye^{-|E_S|^2/8}=k_Y^0,
 \qquad
 \|X'_Y\|=\|\Upsilon_Y\|e^{|D|^2/16+a}.
\]
For the first identity move the left heat factor through the
phase using the Weyl rule. Its resulting squares are
$\frac18|\widetilde E+\delta_1|^2+
 \frac18|\widetilde E-\delta_0|^2$; their linear and constant
terms cancel those in \eqref{eq:diag-two-layer-prime}. Therefore
$Y=G'_{X'_Y}$ and it satisfies the coordinate bound
$2^{|S|/2}\|X'_Y\|$. This bound will also control the approximation
error uniformly in the torus.

For $\epsilon>0$ let $P_\epsilon=e^{-\epsilon|E_S|^2}$ and
\begin{equation}\label{eq:diag-two-layer-epsilon}
 X_{Y,\epsilon}=e^{i(\delta_0+\delta_1)A}
  \exp\left(-\epsilon|\widetilde E+\delta_1|^2
       -\epsilon|\widetilde E-\delta_0|^2
       +\tfrac12\langle\widetilde E,D\rangle+2a\right)\Upsilon_Y.
\end{equation}
The quadratic part is strictly negative, so this too is in
$\mathfrak W(S,\Lambda)$. The same Weyl calculation gives
\begin{align}\label{eq:diag-two-layer-approximation}
 e^{-|E_S|^2/4}X_{Y,\epsilon}e^{-|E_S|^2/4}
                      &=P_\epsilon k_Y^0P_\epsilon,\notag\\
 Y-G_{X_{Y,\epsilon}}&=
               G'_{X'_Y-P_\epsilon X'_YP_\epsilon}.
\end{align}
The second equality follows from the first and the primed
representation, since their continuous kernels agree. Its
coordinate bound is $2^{|S|/2}\xi_Y(\epsilon)$, where
\begin{equation}\label{eq:diag-two-layer-error}
 \begin{split}
 \xi_Y(\epsilon)={}&\|\Upsilon_Y\|
       \sup_{x\in\R^S}e^{-|x|^2/4+\langle x,D\rangle/4+a}\\
 &\quad\cdot(1-e^{-\epsilon|x+\delta_1|^2
                         -\epsilon|x-\delta_0|^2})
       \le C_Y\epsilon.
 \end{split}
\end{equation}
The last inequality uses $1-e^{-u}\le u$ and the boundedness
of a Gaussian times a quadratic polynomial. By contrast,
\begin{equation}\label{eq:diag-two-layer-operator-size}
 \|X_{Y,\epsilon}\|
 =\|\Upsilon_Y\|
   \exp\left(\frac{(1/2-2\epsilon)^2|D|^2}{8\epsilon}
                  +(1/4-\epsilon)(|\delta_0|^2+|\delta_1|^2)\right).
\end{equation}
It can grow exponentially in $1/\epsilon$. This forces the
order of the limits in the next result.

Reflection changes $(\delta_0,\delta_1,\delta_\perp)$ to
$(-\delta_1,-\delta_0,\delta_\perp)$ and
$\Upsilon_Y$ to $\Upsilon_Y^*$. The latter follows by reflecting
\eqref{eq:diag-two-layer-matrix}; $\operatorname{Ins}_c[\omega]$
is even and invariant. Taking adjoints in the two heat identities
then gives
$X'_{\Theta_cY}=(X'_Y)^*$ and
$X_{\Theta_cY,\epsilon}=X_{Y,\epsilon}^*$; the heat operator
is injective with dense range, so these identities determine
the bounded middle operators uniquely. All constructions extend
linearly to two-layer polynomials. Coordinate permutations
give the same assertions at every coordinate link plane.

\begin{proposition}\label{prop:diag-two-layer-principle}
For every $Y\in\mathcal P_1(c)$ at a coordinate link plane,
\begin{equation}\label{eq:diag-two-layer-principle}
 \lim_{t\downarrow0}\limsup_{i\to\infty}
    |\mu_i^t(\Theta_cY)-\overline{\mu_i^t(Y)}|=0.
\end{equation}
The assertion concerns the finite-volume functionals on fermion
polynomials; it does not presume that they converge to fermion
laws.
\end{proposition}

\begin{proof}
For a fixed support, $o_i^t(X)=\mu_i^t(G_X)$ is a unital linear
functional on the local operator algebra with
$|o_i^t(X)|\le\gamma\|X\|$, where
$\gamma=\gamma_{t,i}(k)$. A unital functional this close in norm
to one is nearly Hermitian. If $X_0=X_0^*$ and $\|X_0\|\le1$,
write $o_i^t(X_0)=u+iv$. For each real $s$,
\[
 (v+s)^2\le |o_i^t(X_0+is)|^2
            \le\gamma^2\|X_0+is\|^2
            \le\gamma^2(1+s^2).
\]
For $\gamma>1$ take $s=v/(\gamma^2-1)$; this gives
$|v|\le\sqrt{\gamma^2-1}$. For $\gamma=1$ the same conclusion
follows by letting $s$ tend to either infinity. Decomposing a
general $X$ into its real and imaginary self-adjoint parts yields
\begin{equation}\label{eq:diag-near-hermitian}
 |o_i^t(X^*)-\overline{o_i^t(X)}|
           \le4\sqrt{\gamma^2-1}\|X\|.
\end{equation}
Use \eqref{eq:diag-insertion-reflection} to interpret this as
reflection covariance of the insertion expectation.

Apply \eqref{eq:diag-near-hermitian} to $X_{Y,\epsilon}$ and
use \eqref{eq:diag-two-layer-approximation} for $Y$ and its
reflection. The amplification theorem controls the two errors
and gives
\begin{equation}\label{eq:diag-two-layer-finite-error}
 \begin{split}
 |\mu_i^t(\Theta_cY)-\overline{\mu_i^t(Y)}|
 \le{}&4\sqrt{\gamma^2-1}\|X_{Y,\epsilon}\|\\
      &+2\cdot2^{|S|/2}\gamma\xi_Y(\epsilon).
 \end{split}
\end{equation}
The support box and its coarse scales do not depend on
$\epsilon$. First take $i\to\infty$, and then $t\downarrow0$
with $\epsilon$ fixed. Since $\Delta_i\to0$, the first term
vanishes and the remaining upper bound is
$2\cdot2^{|S|/2}\xi_Y(\epsilon)$. Let $\epsilon\downarrow0$
using \eqref{eq:diag-two-layer-error}. Finite sums and coordinate
transport give the stated result.
\end{proof}

\subsection{Removing fields from a thin layer}\label{subsec:diag-fermion-recursion}

Consider the six thin layers $\mathcal X_{j,a}$ defined above.
Each site has one neighbor in the next thin layer, one in the
preceding one and two in its own layer. The forward matching
bonds have types $2,1,0,1,2,3$ in succession. This follows
directly from the chamber reflections: type zero reverses
the coordinate of modulus $v_3$, types one and two interchange
the adjacent moduli, and type three reverses the coordinate of
modulus $v_1$ and shifts its nearest photon site by the corresponding
unit vector. The bonds of types zero, one and two carry no photon
phase. Type three bonds within thin layers $2,\ldots,5$ carry
spatial phases in plane $j$, while those between layer six and
the next layer one carry temporal phases between $j$ and $j+1$.
These facts control which planes the recursion can read.

For a photon site $y$, let $\mathcal P^\pm(y)$ consist of the
same exponential monomials as $\mathcal P$, with
$d^*\delta=q(\mathfrak q(\Psi)\mp1_y)$. They transform with
gauge phase $e^{\pm iq\lambda(y)}$ and are odd. Removing a
$\bar\psi(x)$ from an invariant monomial leaves an element of
$\mathcal P^+(\pi(x))$; removing a $\psi(x)$ leaves one in
$\mathcal P^-(\pi(x))$. Multiplication by a hopping phase transports
the specified charge to the other endpoint. These assertions
are the identities
\[
 d^*1_{(y,y')}=1_{y'}-1_y,\qquad
 U_A(x,x')=e^{iq\langle1_{(\pi(x),\pi(x'))},A\rangle}.
\]
They show that every term of the following replacements stays
in $\mathcal P$.

Berezin integration by parts, with left derivatives, gives
\begin{align}\label{eq:diag-equations-of-motion}
 Z_T[\Phi(\bar\psi \mathcal D_A)_\alpha(x')G]
       &=-Z_T[\Phi\partial_{\psi_\alpha(x')}G],\notag\\
 Z_T[\Phi(\mathcal D_A\psi)_\alpha(x')G]
       &= Z_T[\Phi\partial_{\bar\psi_\alpha(x')}G].
\end{align}
The appropriate charge of $G$ makes each integrand invariant.
These identities first hold at fixed complex $A$: the
derivatives of $e^{-S_T(A)}$ are respectively
$(\bar\psi \mathcal D_A)_\alpha e^{-S_T(A)}$ and
$-(\mathcal D_A\psi)_\alpha e^{-S_T(A)}$, and the Berezin integral
annihilates a derivative. Integrating against any bounded
bosonic factor $\Phi(F)$ then proves
\eqref{eq:diag-equations-of-motion}. The enlarged local box
must embed so that all generators involved are distinct torus
generators and the phases can be read on the local lift.

Let $x'$ be the neighbor of $x$ in the next thin layer. Its
matching hopping matrix $w(x,x')$ has eigenvalues
$\frac{w_i}{2}(\rho\mp1)$ and is invertible for our fixed
$0<\rho<1$. Solve the equation at $x'$ for the field at $x$.
With the replaced field moved to the left of the monomial,
the two rules are
\begin{equation}\label{eq:diag-field-replacement}
 \begin{split}
 \bar\psi(x)\longmapsto{}&
 \left(M\bar\psi(x')-
   \sum_{x''\sim x',\,x''\ne x}\bar\psi(x'')w(x'',x')U_A(x'',x')
                  +\partial_{\psi(x')}\right)
                     w(x,x')^{-1}U_A(x',x),\\
 \psi(x)\longmapsto{}&U_A(x,x')w(x',x)^{-1}
 \left(M\psi(x')-
   \sum_{x''\sim x',\,x''\ne x}w(x',x'')U_A(x',x'')\psi(x'')
                  -\partial_{\bar\psi(x')}\right).
 \end{split}
\end{equation}
The derivatives act on the monomial to their right. Multiplying
by the inverse matching matrix and its inverse phase in
\eqref{eq:diag-equations-of-motion} proves equality of the
expectations before and after either replacement, for every
bosonic factor $\Phi$. This equality is independent of a choice
of time direction or a symmetry of the torus.

For a thin layer $\mathbb L$ define a linear map $R^+_{\mathbb L}$
by recursion on the number of its fields in an ordered monomial.
Select the first such field, move it to the front with its
Grassmann sign, and apply \eqref{eq:diag-field-replacement}
through its matching neighbor in $\mathbb L^+$. That neighbor
has only one neighbor in $\mathbb L$, the field just removed.
All other fields created by the rule lie in $\mathbb L^+$ or
$\mathbb L^{++}$; its derivatives remove fields in
$\mathbb L^+$. Thus the number in $\mathbb L$ falls by exactly
one at every recursive call. The map is well-defined, preserves
$\mathcal P$, removes all fields in $\mathbb L$, and creates
no new fields outside its next two layers. A downward map
$R^-_{\mathbb L}$ is obtained using the preceding matching
neighbor.

If the initial polynomial is read in a box $B$, the entire
one-layer recursion is read in $B+[-1,1]^3$. The fields being
removed stay in the original box, and their next neighbors
and the other neighbors of those sites have photon projections
equal or nearest to the original projection: among two successive
different bonds at most one has type three. Hence these sites
and every added phase are in the same enlarged box, irrespective
of the number of recursive calls. Induction therefore gives
\begin{equation}\label{eq:diag-thin-layer-identity}
 Z_T[\Phi P]=Z_T[\Phi R^\pm_{\mathbb L}P]
\end{equation}
whenever this enlarged box embeds. The new photon links are
only those of type three bonds with an endpoint in
$\mathbb L^\pm$. The recursion commutes with multiplication
by a curvature character, since it differentiates only fields
and its choices depend only on the monomial.

We first use this recursion to bound an arbitrary polynomial;
here its photon links need not yet be reduced. Move its fields
upwards through all thin layers below $\mathcal X_{j,6}$ and
downwards through those above $\mathcal X_{j+1,1}$. The last
downward replacement can reach $\mathcal X_{j,6}$, and cannot
reach a lower layer. We obtain, on every sufficiently large
torus and with every bounded bosonic factor, an equality
\begin{equation}\label{eq:diag-polynomial-normal-form}
 Z_T[\Phi P]=\sum_a b_a Z_T[\Phi\varphi_{\beta_a}Y_a],
 \qquad Y_a\in\mathcal P_1(c).
\end{equation}
Here the $Y_a$ can be chosen with integer exponents $qj_a$.
Indeed each remaining monomial has integer charge of total
zero on the two adjacent photon planes. In their connected
rectangular support graph choose paths from a fixed vertex
to each charged vertex; the charge-weighted sum of these
paths is an integer cochain $j_a$ with the required divergence.
Then $u_a=\delta_a-qj_a$ is divergence free in the same
integer box and can be written $d^*\beta_a$ there.

The filling must stay inside the support box for the polynomial
bound below. In
$\prod_\mu[a_\mu,a'_\mu]$, set electric components
$\beta_j^{\mathrm{el}}=-\sum_{l=a_1}^j(u_a)_l$ for
$a_1\le j<a'_1$. Subtracting $d^*\beta$ removes all spatial
links below the top plane; divergence zero then removes
the temporal links. On that top plane a second partial-sum
sweep in the second coordinate removes the third-direction
links below its top edge. Divergence zero removes the
second-direction links and leaves only a terminal segment
in direction three. Its divergence equation and vanishing
boundary value make it zero. The two sweeps use only
plaquettes of the box. This proves the asserted filling
and thus the splitting $\delta_a=d^*\beta_a+qj_a$ in
\eqref{eq:diag-polynomial-normal-form}.

Apply the coordinate bound to each $Y_a$, absorb the extra
character and use amplification. For any fixed box $B_0$
there are constants $C_P$, $k_P$ and a volume threshold with
\begin{equation}\label{eq:diag-polynomial-apriori}
 \sup_{\supp\beta\subset B_0}
       |\mu_i^t(\varphi_\beta P)|
             \le C_P\gamma_{t,i}(k_P),\qquad t\ge0.
\end{equation}
One may take
$C_P=\sum_a|b_a|2^{|S_a|/2}\|X'_{Y_a}\|$, with the summand
one for $Y_a=1$. All support boxes in the finite normal form
and $B_0$ fit in one fixed box; its radius determines $k_P$.
The size of $\beta$ does not enter the bound. Reflected
polynomials have the same kind of bound, since the normal
form identity can be reflected on the reflected torus and
the coordinate bounds are preserved by reflection.

\subsection{The contour shift and an outer plane}\label{subsec:diag-contour}

A fermion recursion removes fields from a plane but leaves its
photon exponent. An imaginary shift of links in the next plane
removes that exponent together with the outgoing temporal one.
The all-flux measure is essential for this shift: summation over
exact integer fluxes unfolds the compact link intervals to
real lines on which the contour can move.

Let $u$ be a real cochain supported on finitely many links $S$
of one coordinate plane, and let $P\in\mathcal P$. Choose a
box $B$ reading $P$ and containing $b+[-1,1]^3$ for every
$b\in S$. Put $P^{(u)}(\widehat A)=P(\widehat A-iu)$ by
its entire exponential expansion. Define
\begin{equation}\label{eq:diag-shift-fermion-polynomial}
 \begin{split}
 Q_u=\sum_{b\in S}\sum_{(x,x')\text{ over }b}\big[&
 (e^{qu_b}-1)\bar\psi(x)w(x,x')U_A(x,x')\psi(x')\\
 &+(e^{-qu_b}-1)\bar\psi(x')w(x',x)U_A(x',x)\psi(x)\big],
 \qquad V_u=e^{Q_u}.
 \end{split}
\end{equation}
The bonds are oriented with $b$. Each term is invariant and
even, and the exponential is a finite polynomial by nilpotence.
It reads only those links of $S$ with $u_b\ne0$ and their
overlying bonds. For a plane normal to $e_1$, those fields
are in thin layers $2,\ldots,5$ of its thick layer. The change
in the action is $S_T(A-iu)-S_T(A)=-Q_u$, so
$e^{-S_T(A-iu)}=e^{-S_T(A)}V_u$.

If $\Phi(F)$ is entire and bounded on each tube of bounded
imaginary curvature, then on every torus in which $B$ embeds
we have
\begin{equation}\label{eq:diag-contour-identity}
 Z_T[\Phi P]=e^{\|du\|^2/2}
       Z_T[\Phi(F-idu)\varphi_{du}(F)V_uP^{(u)}].
\end{equation}
The norm here is the lattice norm, equal to the torus norm
because its support embeds. We supply the contour justification.
For each $b$ choose its electric plaquette leaving the plane
in the normal direction. Its only side in $S$ is $b$, with
coefficient $\pm1$. Consequently $d$ is injective on cochains
supported in $S$, and the action
$z\mapsto z+dk$, $k\in\Z^S$, on closed integer fluxes is free.
Choose orbit representatives $z_0$. The compact $A_S$ intervals
and the sum over $k$ become the integral over $\R^S$, since
$(A,z_0+dk)$ represents the same configuration as
$(A+\frac{2\pi}{q}k,z_0)$.

For fixed $z_0$ and the other real links, choose an integer
primitive of $z_0$ in the local box. Adding it to $A$ gives
a local lift affine in the complex variables $\zeta=A_S$.
The Berezin coefficients of $e^{-S_T(A)}$ are polynomials
in the link phases, so the full integrand extends to an
entire function $f(\zeta)$. On the strip
$|\Im\zeta_b|\le\|u\|_\infty$ the selected plaquettes give
\begin{equation}\label{eq:diag-contour-strip}
 |f(\zeta)|\le C
          \prod_{b\in S}e^{-(\Re\zeta_b-a_b)^2/2}.
\end{equation}
Indeed each of their real curvatures is $\pm\Re\zeta_b+a'_b$,
all imaginary curvatures are bounded in this finite strip,
the factor $\Phi$ is bounded there, and each remaining
exponential phase is bounded there. The constant can be
chosen independent of the remaining real links and the
orbit representative; only the real Gaussian centers vary.

Move one variable at a time along the rectangle with real
endpoints $\pm R$ and imaginary displacement $-iu_b$. The
vertical sides tend to zero by \eqref{eq:diag-contour-strip}.
Fubini then gives $\int f(\xi)\,\dd\xi=
\int f(\xi-iu)\,\dd\xi$. At the shifted arguments the
Maxwell exponent becomes
\[
 -\tfrac12\langle F-idu,F-idu\rangle
       =-\tfrac12\|F\|^2+i\langle du,F\rangle+\tfrac12\|du\|^2.
\]
The change in the fermion action is the polynomial $V_u$,
and the observable becomes $P^{(u)}$. Both the initial
integrand and the resulting real integrand are absolutely
integrable by their local bounds. Hence we may sum the
identities over $z_0$ and integrate the remaining links,
which proves \eqref{eq:diag-contour-identity}.

We now eliminate the lowest plane $j$ read by $P$. This means
that its fields have photon projection height at least $j$
and its links have no endpoint below plane $j$. Apply the
six upward thin-layer maps of this thick layer to obtain
$P^\uparrow=\sum_a b_a e^{i\delta'_a\widehat A}\Psi_a$.
It reads no fields of thick layer $j$ and no links below it;
the only new links are spatial links in plane $j$ and temporal
links from $j$ to $j+1$. At each term, gauge invariance in
plane $j$ therefore says
\begin{equation}\label{eq:diag-lowest-plane-charge}
 (\delta'_a)^\perp_j=d_s^*(\delta'_a)_j.
\end{equation}
For this term let $u_a$ be supported on plane $j+1$ with
$(u_a)_{j+1}=(\delta'_a)_j$. Define
\begin{equation}\label{eq:diag-outer-polynomial}
 P_a=b_a e^{\|du_a\|^2/2+\langle\delta'_a,u_a\rangle}
       e^{i\langle\delta'_a+d^*du_a,\widehat A\rangle}
                                    V_{u_a}\Psi_a.
\end{equation}
Then \eqref{eq:diag-thin-layer-identity} and
\eqref{eq:diag-contour-identity} give the exact relation
\begin{equation}\label{eq:diag-outer-identity}
 Z_T[\Phi P]=\sum_a Z_T[\Phi(F-idu_a)P_a].
\end{equation}
The construction is on the infinite lattice; the same relation
holds for every $\Phi$ entire in the plaquette variables and
bounded on each tube $|\Im F_\sigma|\le r$ for all $\sigma$,
$r>0$, and every torus containing the fixed enlarged support box.

Here is the cancellation of the eliminated exponent. Write
$v=(\delta'_a)_j$. From the slice formulas
$(du)^{\mathrm{el}}_l=u_{l+1}-u_l-d_su_l^\perp$ and
$(d^*\beta)_l=\beta^{\mathrm{el}}_{l-1}
-\beta^{\mathrm{el}}_l+d_s^*\beta^{\mathrm{mag}}_l$, the only
nonzero components of $d^*du_a$ are
\begin{align}\label{eq:diag-outer-cancellation}
 (d^*du_a)_j&=-v,&(d^*du_a)^\perp_j&=-d_s^*v,\notag\\
 (d^*du_a)_{j+1}&=(2+d_s^*d_s)v,
       &(d^*du_a)^\perp_{j+1}&=d_s^*v,\notag\\
 (d^*du_a)_{j+2}&=-v.
\end{align}
Thus the spatial exponent in plane $j$ vanishes, and
\eqref{eq:diag-lowest-plane-charge} cancels its outgoing
temporal exponent as well. The polynomial $V_{u_a}$ reads
only thick layer $j+1$. Hence $P_a$ has lowest plane at
least $j+1$ and highest plane at most the maximum of $j+2$
and the previous highest plane. The six field recursions
enlarge the box by six, and the shifted links with their
plaquettes require two more. The whole step is read in
$B+[-8,8]^3$. Elimination from the highest plane is obtained
by the downward maps and a shift in plane $j-1$. The corresponding
incoming temporal exponent is
$-d_s^*(\delta'_a)_j$, and the same two components cancel.
The highest plane falls by one, while the lowest can fall
only to $j-2$.

The source of error in the tilted law is precisely the shifted
bosonic factor. For a finite cochain $\eta$, only finitely
many copies in the array $V_i$ meet its support. Their change
is the fixed trigonometric polynomial
\begin{equation}\label{eq:diag-array-shift}
 H_\eta(F)=\sum_{\kappa,b}a_b
       (e^{\langle(g_\kappa^{-1})^*\beta_b,\eta\rangle}-1)
                        \varphi_{(g_\kappa^{-1})^*\beta_b}(F),
\end{equation}
where $O=\sum_ba_b\varphi_{\beta_b}$ and only contributing
cells occur. At most $|\supp\eta|$ cells contribute; each
lies in the support box of $\eta$ enlarged by $h$. Its
coefficient norm has the bound
\begin{equation}\label{eq:diag-array-shift-size}
 |H_\eta|_1\le C(\eta):=|\supp\eta|\,|O|_1
        (e^{\max_b\|\beta_b\|_1\|\eta\|_\infty}-1).
\end{equation}
If that enlarged box embeds, no two contributing cells are
period equivalent and the same polynomial is the torus
difference $V_i(F-i\eta)-V_i(F)$. Expand its exponential
in characters:
\begin{equation}\label{eq:diag-shifted-tilt}
 e^{tV_i(F-i\eta)}=e^{tV_i(F)}
        \left(1+\sum_\alpha b_\alpha(t)\varphi_{\beta_\alpha}(F)\right),
 \qquad\sum_\alpha|b_\alpha(t)|\le e^{tC(\eta)}-1.
\end{equation}
The coefficient norm is submultiplicative, so the exponential
series proves absolute uniform convergence. All character
supports lie in one fixed enlarged box, independent of $i$
and $t$; sums of character exponents do not enlarge this support.

Set $P^\flat=\sum_aP_a$ in \eqref{eq:diag-outer-identity}.
Insert \eqref{eq:diag-shifted-tilt} with $\eta=du_a$ and divide
by the positive denominator. Termwise integration is valid:
at a fixed torus the functional on bounded bosonic factors
multiplying $P_a$ is continuous for their supremum norm, by
the absolute integrability of its Berezin coefficients.
We get
\begin{equation}\label{eq:diag-outer-error}
 |\mu_i^t(P)-\mu_i^t(P^\flat)|
 \le(e^{Ct}-1)\sum_a
       \sup_{\supp\beta\subset B'}|\mu_i^t(\varphi_\beta P_a)|,
\end{equation}
where $C=\max_aC(du_a)$ and $B'$ is a fixed box enlarged
by $h$. The a priori bound
\eqref{eq:diag-polynomial-apriori} makes this $O(t)$ uniformly
for all large $i$ and $0\le t\le1$: each of the finitely many
amplification factors is bounded once $\Delta_i\le1$, and
$e^{Ct}-1\le(e^C-1)t$ on this interval.

The same estimate applies after transporting the identity.
This needs care because $D_i$ need not be invariant under
a coordinate permutation. Reflect the equality on the
reflected torus and then reflect it back, using functional
covariance. The conjugations change a shift $\eta$ to
$-\theta^*\eta$. Indeed, if
$\Phi^\theta(F)=\overline{\Phi(\theta^*\overline F)}$, then
$(\Phi^\theta(\,\cdot-i\eta))^\theta=
\Phi(\,\cdot+i\theta^*\eta)$. Thus
\eqref{eq:diag-outer-identity} becomes an exact identity
for $\Theta P$ with shifts $-\theta^*du_a$ and polynomials
$\Theta P_a$ on the original torus. Iterating reflections
transports it by any $g$, with shifts
$\det(\vec g)(g^{-1})^*du_a$, where $\vec g$ is the linear
part of $g$. Its support box is transported at the same time. Apply \eqref{eq:diag-shifted-tilt} and
\eqref{eq:diag-polynomial-apriori} on $D_i$ itself to obtain
the transported $O(t)$ estimate. No invariance of $D_i$
or of its tilted array under $g$ is used.

\begin{proposition}\label{thm:diag-polynomial-reduction}
For each $P\in\mathcal P$ and coordinate link plane $c$ there
is $Y_P\in\mathcal P_1(c)$, independent of the torus and tilt,
such that, for every fixed lattice isometry $g$,
\begin{equation}\label{eq:diag-polynomial-reduction}
 |\mu_i^t(P^g)-\mu_i^t((Y_P)^g)|\le C_{P,g}t
 \quad(0\le t\le1, i\ge i_0(P,g)).
\end{equation}
At a plane normal to either other coordinate the corresponding
reduction and its reflected estimate hold as well.
\end{proposition}

\begin{proof}
For $c=j_0+\tfrac12$, eliminate lowest planes until the
lowest is at least $j_0$. Every step raises it by at least
one, so this takes finitely many steps. Then eliminate
highest planes until the highest is at most $j_0+1$.
During this phase an eliminated plane $j\ge j_0+2$ can
create links only as low as $j-2\ge j_0$, preserving the
first reduction. The resulting polynomial reads just
the two thick layers and their connecting links. Apply
the upward recursions in thin layers $1,\ldots,5$ of
the lower thick layer, followed by the downward recursions
in thin layers $6,\ldots,2$ of the upper one. New fields
remain at or above the lower layer six and at or below
the upper layer one. The new type three phases are
spatial phases in these two planes or the temporal
phases between them. Thus the result $Y_P$ is in
$\mathcal P_1(c)$.

There are finitely many outer-plane steps, each with
\eqref{eq:diag-outer-error} and its transported version;
the final ten thin-layer recursions are exact with the
unchanged bosonic factor and contribute no error. Take
one volume threshold for all their fixed support boxes
and a priori bounds, and sum the finite $O(t)$ errors.
This proves \eqref{eq:diag-polynomial-reduction}.
For another direction first apply the construction to
$P^{g_\mu^{-1}}$ and then transport its exact identities
by the cyclic coordinate permutation $g_\mu$. The general
transported estimate, evaluated on $D_i$, proves the
claim there and its reflected version.
\end{proof}

